Un treuil électrique de \qty{\v{P}}{W} soulève une caisse de \qty{\v{m}}{kg} à vitesse constante sur une hauteur de \qty{\v{h}}{m} en \qty{\v{t}}{s}. On utilise $g = \qty{9.8}{N/kg}$.

\begin{enumerate}[label=\alph*)]
    \item Quelle quantité d'énergie électrique le moteur consomme-t-il pendant la montée ?
    \item Quelle est l'énergie utile, c'est-à-dire le gain d'énergie potentielle de la caisse ?
    \item Calculez le rendement énergétique du treuil.
    \item Quelle quantité d'énergie n'est pas utile ? Sous quelle forme se retrouve-t-elle surtout ?
\end{enumerate}
